% Generated draft: requires mathematical and source review.
\documentclass[11pt]{article}
\usepackage[margin=25mm]{geometry}
\usepackage{fontspec}
\setmainfont{Latin Modern Roman}
\usepackage{amsmath,amssymb,mathtools,unicode-math}
\setmathfont{Latin Modern Math}
\usepackage{enumitem,needspace,xurl,hyperref}
\setlength{\parindent}{0pt}
\setlength{\parskip}{5pt}
\setlength{\emergencystretch}{3em}
\newcommand{\N}{\mathbb{N}}
\newcommand{\Z}{\mathbb{Z}}
\newcommand{\Q}{\mathbb{Q}}
\newcommand{\R}{\mathbb{R}}
\newcommand{\C}{\mathbb{C}}
\newcommand{\F}{\mathbb{F}}
\newcommand{\A}{\mathbb{A}}
\newcommand{\PP}{\mathbb{P}}
\newcommand{\E}{\mathbb{E}}
\newcommand{\eps}{\varepsilon}
\newcommand{\dd}{\,\mathrm{d}}
\begin{document}
\section*{UnsolvedMath ID: 20000324}
\textbf{Source problem code:} catalogue entry
\subsection{Definitions and mathematical statement}
A stable surface is, in the intended moduli convention, a projective surface with semi-log-canonical singularities and ample \(\Q\)-Cartier canonical divisor. Semi-log-canonicity includes the appropriate reduced \(S_2\), codimension-one and normalization/conductor log-canonicity conditions. A family must satisfy a specified \(\Q\)-Gorenstein/reflexive-base-change condition. The source, attributed to Hassett, asks for the valuative criterion for properness of their positive-characteristic moduli stack, without fixing the moduli functor, numerical component or allowable fields.

For a fixed bounded numerical component, a valuative formulation asks that a stable surface over the fraction field of a discrete valuation ring extend, after an allowed finite extension of valuation rings, to a stable family, and that an isomorphism of generic fibers of two stable families extend uniquely over the valuation ring. The precise stack-theoretic existence and uniqueness conventions must be fixed; finite base change in stable reduction is not silently forbidden. Unbounded all-volume stacks need not be quasi-compact, so finite-type properness and componentwise valuative properness must be distinguished.

The attached aid describes the unqualified 2004 question as only partially solved: fixed-volume stable reduction is reported in equal characteristic sufficiently large relative to the volume/data, and uniform properness over \(\Z[1/30]\) is reported conditional on locally stable reduction. No complete publication citations or thresholds are supplied.

Its numerical lemma assumes an explicit log-canonical total model, \(\Q\)-Cartier adjunction and a DCC volume-gap input. Write the special fiber as \(\sum_i m_iD_i\), with positive integers \(m_i\), let \(v>0\) be generic volume, and suppose the relevant adjunction volumes \(v_i\) satisfy the weighted comparison \(v\geq\sum_i m_iv_i\) and \(v_i\geq\delta>0\). In that supplied conditional setting it reports
\[
m_i\leq B:=\left\lfloor\frac v\delta\right\rfloor,\qquad
e:=\operatorname{lcm}_i(m_i)\mid
N_{\mathrm{univ}}:=\operatorname{lcm}(1,\ldots,B),
\]
with \(B\geq1\) whenever these nonempty-fiber hypotheses hold. It calls \(e\) the actual multiplicity-killing Kummer index. The aid reports that \(N_{\mathrm{univ}}\) is the least positive integer forced to be divisible by every multiplicity allowed solely by \(1\leq m\leq B\). This is arithmetic minimality from that bound, not a claim that every allowed multiplicity occurs geometrically. If \(p>B\), both indices are prime to \(p\), and the resulting Kummer base change is reported tame and separable. None of this supplies the missing total model, volume gap or complete stable reduction in every characteristic.
\subsection{Short English statement}
Prove stable-surface reduction and uniqueness in positive characteristic for a specified moduli component. The attached aid reports partial properness results and, conditional on a suitable total model and volume gap, bounds special-fiber multiplicities and gives the least universal arithmetic Kummer index needed to divide all of them.
\subsection{Sources}
\begin{itemize}
\item Local UnsolvedMath registry, permanent ID \(20000324\), exact public code \texttt{catalogue entry}; canonical location \texttt{aim-algebraic-geometry-notes.json}, notes[323].
\item AIM workshop ``Compact moduli spaces and birational geometry,'' extracted text ``Problem 1. 5. (Hassett)''; supplied URL \url{https://aimath.org/WWN/birational/birational.pdf}. The problem number is unresolved. The PDF was not consulted here; the original literature field is empty.
\item Attached research attempt 1, ``Stable-surface properness and an optimal universal tame Kummer index,'' describing the historical question as from 2004 and reporting fixed-volume/equal-characteristic and conditional \(\Z[1/30]\) results without full bibliographic citations.
\end{itemize}
\end{document}
